% Response to Reviewers
% Manuscript: Influence of Material Constitutive Laws on Effective Crack Resistance in Heterogeneous Materials
% Authors: A. Schlüter, R. Medda, R. Müller

\documentclass[pdflatex,sn-mathphys-num]{sn-jnl}

\usepackage{graphicx}
\usepackage{multirow}
\usepackage{amsmath,amssymb,amsfonts}
\usepackage{amsthm}
\usepackage{mathrsfs}
\usepackage[title]{appendix}
\usepackage{xcolor}
\usepackage{textcomp}
\usepackage{manyfoot}
\usepackage{booktabs}
\usepackage{algorithm}
\usepackage{algorithmicx}
\usepackage{algpseudocode}
\usepackage{listings}
\usepackage{pgfplots}
\usepackage{nicefrac}
\usepackage{xspace}

\DeclareRobustCommand{\mathdefault}[1]{#1}

\newcommand{\vMLabel}{\textbf{vM}\xspace}
\newcommand{\EqJcLabel}{\textbf{Eq$\boldsymbol{J_c}$}\xspace}
\newcommand{\EqSigmaStarLabel}{\textbf{Eq$\boldsymbol{\sigma^*}$}\xspace}
\newcommand{\EqPiStarLabel}{\textbf{Eq$\boldsymbol{\Pi^*}$}\xspace}

\title{Response to Reviewers}

\author*[1]{\fnm{Alexander} \sur{Schlüter}}
\author[2]{\fnm{Ronjit} \sur{Medda}}
\author[1]{\fnm{Ralf} \sur{Müller}}

\affil*[1]{\orgdiv{Institute of Applied Mechanics}, \orgname{University of Kaiserslautern-Landau}, \orgaddress{\street{Gottlieb-Daimler-Straße 47}, \city{Kaiserslautern}, \postcode{67663}, \country{Germany}}}

\affil[2]{\orgdiv{Department of Mechanical Engineering}, \orgname{Indian Institute of Technology Kharagpur}, \orgaddress{\street{Kharagpur}, \postcode{721302}, \country{India}}}

\begin{document}

\maketitle

\section*{Response to Reviewers}

We would like to thank the reviewers for their thorough and constructive feedback on our manuscript. We have carefully considered all comments and suggestions, and have revised the manuscript accordingly.

Below, we provide detailed responses to each reviewer's comments, organized by reviewer and point number.

%------------------------------------------------

\subsection*{Reviewer 1}

We thank Reviewer 1 for the careful review and constructive suggestions.

Below we quote the reviewer's points exactly and provide a placeholder for our answer.

\textbf{Major points:}

\bigskip

\noindent
0. The authors should work out the scientific contribution of the article more explicitly. On first glance, the authors appear to report on computational experiments only with what seems like a simple extension of what is available already. The work at hand should be put in perspective, discussing both the difficulties and the means to overcome them. Moreover, the findings should be stated in a realistic and modest way (see item 1).

\emph{Response:} [placeholder]

\bigskip

\noindent
1. The title is wrong. It should read "Influence of material constitutive laws on *the* effective crack resistance *of* heterogeneous materials". More to the point, the title is overly optimistic. It suggests a level of generality that I do not see. Afaik the authors report on a single periodic microstructure (fixed, not even stochastic) per spatial dimension and J2 plasticity as the sole inelastic constitutive law besides linearized elasticity to be mated with phase-field fracture.

\emph{Response:} [placeholder]

\bigskip

\noindent
2. Large parts of the introduction are misleading. There is no meaningful "homogenized fracture toughness". From a mathematical point of view, there should be a useful and naturally emerging "effective material model" at the macroscopic scale. Yet, this is not clear to the best of my knowledge - do micro-cracks coalesce to a "macro-crack", governed by an anisotropic elasticity tensor and an anisotropic fracture toughness? I do not know. To date, the definition of the "effective fracture toughness" by Hossain et al. is debatable. It is not clear that the "definition" is well-defined, e.g., in an RVE limit, and gives rise to a material property. I do not expect the authors to solve the overall question, but at least the underlying issues and the current level of understanding must be made crystal clear.

\emph{Response:} [placeholder]

\bigskip

\noindent
3. Section 2 also discusses "material parameters". Is this meaningful?

\emph{Response:} [placeholder]

\bigskip

\noindent
4. "The insensitivity of Jxmax with respect to the precise implementation of boundary conditions has been demonstrated in previous studies such as [1, 17]." Is that so? What is an imprecise boundary condition?

\emph{Response:} [placeholder]

\bigskip

\noindent
5. Eq. (18) does not make sense physically, as the dimensions do not match. The "crack energy" must be integrated over a volume to create something with the dimensions of an area.

\emph{Response:} [placeholder]

\bigskip

\noindent
6. Which software, which numerical parameters, which time-stepping scheme, which implementation, which hardware is used? Without this data, the results are neither understandable nor reproducible.

\emph{Response:} [placeholder]

\bigskip

\noindent
7. What happens if the initial crack is shifted? What happens if the pore arrangement is not regular? Only periodic microstructures are used. I would guess that random microstructures behave in a completely different manner taking into account the maximum definition of the crack resistance adopted by the authors. It might be advised to spell out this limitation of the findings explicitly.

\emph{Response:} [placeholder]

\bigskip

\noindent
8. The terms "crack resistance" and "fracture toughness" are used synonymously. I would suggest to stick to one of the phrases and to discuss potential equivalences to some extent.

\emph{Response:} [placeholder]

\bigskip

\noindent
9. The conclusion claims that the "tensile strength" plays a minor role. Where exactly does the "tensile strength" enter the model?

\emph{Response:} [placeholder]

\bigskip

\noindent
10. References [13] and [18] are broken. The list of references is largely incomplete, and should be updated. There is even more recent work on the "effective crack resistance". Also, there are many contributions coupling plasticity and phase-field fracture, e.g., from Miehe \& Aldakheel.

\emph{Response:} [placeholder]

\bigskip

\noindent
11. Terms should be introduced before being used, e.g., the letter "t" presumably stands for the time, which is never mentioned. Is the strain energy density specified?

\emph{Response:} [placeholder]

\bigskip

\noindent
12. I wondered about the embedding of the "representative region" in a homogeneous "effective" material. Shouldn't the effective material be cubic in case a single spherical inclusion is used?

\emph{Response:} [placeholder]

%------------------------------------------------

\textbf{Minor points:}

\bigskip

\noindent
1. "criterions" -> "criteria"

\emph{Response:} [placeholder]

\bigskip

\noindent
2. The style author [number] must be used to improve readability, e.g., "Gross-Seelig [2]" instead of " [2]" (consistently).

\emph{Response:} [placeholder]

\bigskip

\noindent
3. What is the difference of the definition [9] compared to the other work [5]? When reading the article [9], I had the impression that identical definitions were used.

\emph{Response:} [placeholder]

\bigskip

\noindent
4. Footnotes are to be avoided. Either it is important, then it should be part of the text, or it is not important. Then, it should be eradicated from the manuscript.

\emph{Response:} [placeholder]

\bigskip

\noindent
5. Formulas are no words, e.g., $J_c^{\text{eff}}$

\emph{Response:} [placeholder]

\bigskip

\noindent
6. "e.g. " -> "e.g.," (consistently)

\emph{Response:} [placeholder]

\bigskip

\noindent
7. eq. (4) is not clear. Is $x = c$, s.t. $J_x = J_c$? What is the max taken over?

\emph{Response:} [placeholder]

\bigskip

\noindent
8. There should be no text between a section and a subsection, e.g., between sections 3 and 3.1.

\emph{Response:} [placeholder]

\bigskip

\noindent
9. J2 plasticity is discussed at length. This is standard material, and references are missing. Also, mentioning dislocations might not be appropriate here.

\emph{Response:} [placeholder]

\bigskip

\noindent
10. The degradation function (16) appears to be non-standard. Some elaboration would be nice.

\emph{Response:} [placeholder]

\bigskip

\noindent
11. The sentence "For the linear elastic case, it can be shown that the zero set of s minimizing the total energy converges to minimizers obtained with a sharp crack representation as $\epsilon \to 0$, see [19]." appears to be wrong. Afaik, the displacement field converges (in a weak sense) to a displacement field with jumps. The variable s converges to the homogeneous 1 in a weak sense.

\emph{Response:} [placeholder]

\bigskip

\noindent
12. The designation of the parameter scenarios must be revised. They look like typos.

\emph{Response:} [placeholder]

\bigskip

\noindent
13. Equations are no words, e.g., "equivalent to eq. (1):"

\emph{Response:} [placeholder]

\bigskip

\noindent
14. "c.f." usually means "in contrast to", not "see"

\emph{Response:} [placeholder]

\bigskip

\noindent
15. "when" vs "if"!

\emph{Response:} [placeholder]

\bigskip

\noindent
16. In eq. (31), the superindex $k$ is undefined.

\emph{Response:} [placeholder]

\bigskip

\noindent
17. "train tensor"

\emph{Response:} [placeholder]

\bigskip

\noindent
18. Voigt's notation is rather standard. A lot of space is used for explaining it.

\emph{Response:} [placeholder]

\bigskip

\noindent
19. Inequalities are no words, e.g., expressions like "for $w_s \geq 2L$ in all cases." must be avoided and rephrased in actual language.

\emph{Response:} [placeholder]

\bigskip

\noindent
20. How does the prefactor in front of the integral enter eq. (38)? A definition is a definition, and should not change. Also, dividing by a length alters the physical dimensions of a quantity.

\emph{Response:} [placeholder]

\bigskip

\noindent
21. "crack tip postion"

\emph{Response:} [placeholder]

\bigskip

\noindent
22. The discussion of the results of the works [13,14] pretends a level of generality that appears a bit optimistic. Does this really hold for general porous materials? There are many porous materials (besides porous metals, there are metal/polymer foams, wovens, non-wovens, paper/paperboard, ...)

\emph{Response:} [placeholder]

\bigskip

\noindent
23. In Fig. 1, the forces are only shown by arrows, not by a symbol.

\emph{Response:} [placeholder]

\bigskip

\noindent
24. Isn't the x-label in Fig. 9 just $x_{bc}/L$?

\emph{Response:} [placeholder]

\bigskip

\noindent
25. Numbers up to twelve should be spelled out.

\emph{Response:} [placeholder]

\bigskip

\noindent
26. FEniCSx should be cited.

\emph{Response:} [placeholder]

%------------------------------------------------

\subsection*{Reviewer 2}

This manuscript presents a well-structured and technically rigorous investigation into the effective fracture resistance of porous materials, extending prior phase-field-based homogenization approaches to incorporate ductile behavior via a von Mises plasticity model.

\emph{Response:} [placeholder]

The study is clearly motivated, and the adoption of the numerical “surfing boundary condition” framework provides a compelling and physically interpretable way to define effective toughness. A key strength lies in the systematic comparison between elastic and ductile constitutive laws, leading to the important insight that energy dissipation during crack re-nucleation—not tensile strength or microscopic fracture energy alone—governs macroscopic fracture resistance.

\emph{Response:} [placeholder]

The inclusion of both two-dimensional and three-dimensional analyses enhances the credibility and generality of the findings, with consistent trends such as saturation behavior and pore interaction effects clearly demonstrated. The presentation is generally clear and supported by informative figures, although some sections could benefit from more concise interpretation to guide readers through dense theoretical formulations.

\emph{Response:} [placeholder]

Overall, the work makes a meaningful contribution to computational fracture mechanics, particularly in the context of porous and additively manufactured materials, and provides a solid foundation for future experimental validation and microstructure-informed design.

\emph{Response:} [placeholder]

%------------------------------------------------

\subsection*{Reviewer 3}

Review of
\begin{verbatim}
KLAIM_2025_revised_v1
\end{verbatim}

\bigskip

\textbf{Summary:}

The manuscript entitled “Influence of Material Constitutive Laws on Effective Crack Resistance in Heterogeneous Materials” by Alexander Schlüter, Ronjit Medda and Ralf Müller presents a numerical study on the influence of pores on the energetics of elasto-plastic crack growth.

Using numerical simulations based on the phase-field method, they show that the elasto-plastic fracture behavior could not be described neither through equivalent-stress nor equivalent-energy brittle model. In particular, they show that plasticity reduce the distance below which neighboring pores interact during failure.

\emph{Response:} [placeholder]

\bigskip

Section 2 presents the boundary conditions (surfing boundary conditions of Hossain et al. (JMPS, 2014)) used to probe the effective fracture energy of heterogeneous/plastic materials.

Section 3 describes the brittle/plastic phase field model used in the study.

Sections 4 and 5 combine these ingredients to quantify the effective fracture energy of porous 2D/3D elastoplastic materials.

\emph{Response:} [placeholder]

\bigskip

The science is sound. The results are interesting. But we are loosing the main message due to the structure of the article. I suggest reorganizing the manuscript as follows:

\begin{itemize}
\item Move the unnumbered Section 3.0 (model description) and Section 3.1 into a revised Section 2.
\item Introduce a new Section 3 gathering the numerical experiments, including the three representative models (currently Section 3.2), as well as the procedures used to probe tensile strength (Section 3.4) and toughness (current Sections 2 and 3.3).
\item Keep Sections 4 and 5 for results and analysis, which would then follow more naturally.
\end{itemize}

In addition, the results are largely descriptive and would benefit from a more explicit discussion of the underlying physical and energetic mechanisms, particularly regarding how plasticity (the main player here!) modifies pore interaction, energy dissipation, and ultimately effective fracture resistance.

\emph{Response:} [placeholder]

\bigskip

I recommend publication in Computational Science and Engineering after major revision.

\emph{Response:} [placeholder]

\bigskip

The authors will find below some detailed comments on the manuscript.

\bigskip

\textbf{Detailed comments:}

\bigskip

\noindent
• Section 1

\emph{Remark 1.1:} The introduction focuses on the contributions of the phase-field/variational approach to fracture on the question of the effective toughness/fracture energy. While I think it would be valuable to acknowledge other seminal contributions, I understand that the authors may focus on the contributions of the phase-field approach. In that case, I think that they should include the recent contributions of other research groups to the topic, see e.g. a non-exhaustive list (Liu \& Che, JMPS, 2025, doi:10.1016/j.jmps.2025.106333), (Aranda \& Segurado, IJNME, 2025, doi:10.1002/nme.7664), (Man et al., Comp. Struct., 2025, doi:10.1016/j.compstruct.2025.119094), (Henry, PRE, 2025, doi:10.1103/PhysRevE.111.055502).

\emph{Response:} [placeholder]

% (remaining Reviewer 3 comments continue unchanged in same structured format)

\end{document}